\documentclass[11pt]{article}

\usepackage[utf8]{inputenc}
\usepackage[T1]{fontenc}
\usepackage{lmodern}
\usepackage[margin=1in]{geometry}
\usepackage{booktabs}
\usepackage{array}
\usepackage{amsmath}
\usepackage{graphicx}
\usepackage[hidelinks]{hyperref}
\usepackage[round]{natbib}
\usepackage{setspace}
\usepackage{titlesec}

\setlength{\parskip}{0.4em}
\setlength{\parindent}{0pt}

\title{Segovia: a chunked, GIL-released, bounded-memory streaming engine
for near-real-time multichannel electrophysiology preprocessing}

\author{Felipe Carvajal Brown\\
\small Independent Researcher\\
\small \texttt{fcarvajalbrown@gmail.com} \quad ORCID: 0000-0002-8300-7587}

\date{}

\begin{document}

\maketitle

\begin{abstract}
High-density extracellular electrophysiology probes (Neuropixels) acquire 384--768 channels at
30\,kHz, producing about 22\,MB/s per probe. Preprocessing this stream for near-real-time use, such
as closed-loop stimulation or online brain--machine-interface decoding, means each incoming chunk of
samples must be processed inside its acquisition period while peak memory stays bounded and
independent of how long the recording runs. Batch tools driven one chunk at a time re-read the
filter-margin neighbourhood on every call and give no analytical memory bound. We describe Segovia,
an open-source Rust library (AGPL-3.0) with Python bindings (PyO3) built for this regime. Its core
abstraction is the \texttt{ChunkSource} trait: an iterator over fixed-size \texttt{i16} chunk buffers
consumed by a Rayon-parallelized bandpass-filter, common-median-reference, and ZCA-whitening chain.
The Python Global Interpreter Lock is released for the duration of each Rust computation, and peak
memory is bounded analytically by chunk size rather than recording length. We evaluate Segovia on a
real International Brain Laboratory Neuropixels AP-band recording (385 channels, mtscomp-compressed)
and on a built-in streaming synthetic simulator (\texttt{SyntheticEphysReader}), using a
replay-at-acquisition-rate harness that measures per-chunk latency and deadline adherence without
hardware. Streaming the full 55.8-minute recording at a 300\,ms chunk budget, Segovia meets its
real-time deadline on 99.7\% of chunks at 0.21\,GB peak memory, against 94.7\% at 0.41\,GB for an
equivalent SpikeInterface online-streaming configuration; the larger separations are in peak memory
(2$\times$), maximum latency (334 vs 932\,ms), and jitter (39 vs 60\,ms). In the cold-start window the
adherence gap is wider (100\% vs 69.5\% at the same budget) because SpikeInterface pays its per-chunk
setup cost most heavily before the run warms up. Segovia is available at
\url{https://github.com/fcarvajalbrown/Segovia} and via \texttt{pip install segovia}.
\end{abstract}

\section{Introduction}

Neuropixels probes \citep{Jun2017} have made large-scale, high-density extracellular
electrophysiology routine. A single probe acquires 384 channels at 30\,kHz for about 22\,MB/s, and
multi-probe rigs multiply that by four to eight. A one-to-two-hour session produces tens of gigabytes
of raw data. Before spike sorting can begin, this stream passes through a preprocessing chain:
bandpass filtering to isolate action-potential frequencies (300--6000\,Hz), common-median referencing
(CMR) to suppress common-mode noise, and whitening to decorrelate channels. In the standard offline
pipeline (SpikeInterface \citep{Buccino2020}, MountainSort), this step runs after recording finishes
and is tuned for throughput on a fixed dataset.

Near-real-time applications work under different constraints. Closed-loop optogenetic stimulation must
detect a neural event and respond within tens to hundreds of milliseconds, inside the same chunk of
samples the event occurred in. Online brain--machine-interface decoders must emit decoded signals with
bounded latency. Hardware-in-the-loop experiments must synchronize preprocessing output with external
devices at acquisition rate. These uses need two guarantees that batch pipelines do not offer:

\begin{enumerate}
  \item \textbf{Real-time deadline adherence.} Each chunk of $N$ samples must be preprocessed within
  $N / f_s$ seconds, its chunk period. Missing this deadline stalls the application.
  \item \textbf{Bounded, file-size-independent peak memory.} A closed-loop system may run for hours or
  days, so it cannot exhaust RAM as the recording grows.
\end{enumerate}

Drive a batch tool such as SpikeInterface one chunk at a time through sequential \texttt{get\_traces}
calls, the natural online adaptation, and each call re-reads and re-decodes the filter-margin
neighbourhood (the look-ahead and look-behind samples the filter needs). Whitening-matrix estimation
may also be repeated or anchored to cached state in ways that were never meant for streaming. Memory
then grows with the number of processed chunks or with the mapping of the whole recording, not with
chunk size alone. Segovia is a purpose-built streaming engine that meets both constraints by
construction.

\section{Related work}

\textbf{Batch preprocessing frameworks.} SpikeInterface \citep{Buccino2020} gives a unified Python API
for spike sorting that includes the standard preprocessing chain. Its \texttt{ChunkRecordingExecutor}
spreads processing across chunks with a thread or process pool, tuned for bulk throughput. It does not
target per-chunk latency or an analytical memory bound. MountainSort \citep{Chung2017} follows a
similar offline model. Kilosort4 \citep{Pachitariu2024}, the current state-of-the-art GPU spike
sorter, works on already-preprocessed recordings rather than raw streams. None of these target the
online, one-chunk-at-a-time regime.

\textbf{Streaming and closed-loop frameworks.} improv \citep{Draelos2025} is a Python framework for
closed-loop calcium-imaging experiments, benchmarked by replaying prerecorded data at acquisition rate
rather than against live hardware. BRAND \citep{Ali2024} is a modular platform for closed-loop
experiments with deep-network models, reaching sub-600\,\textmu s latency over asynchronous
Redis-based messaging. RT-Sort \citep{vanderMolen2024} does real-time spike detection and sorting with
7.5\,ms latency on recorded ground-truth datasets. Together these establish replay-at-acquisition-rate
as an accepted way to evaluate near-real-time systems without a live rig \citep{Hoefler2015}.

\textbf{Simulators.} MEArec \citep{BuccinoMEArec2020} generates synthetic extracellular recordings
with ground-truth spike labels from biophysical neuron models. Segovia's built-in simulator borrows
the MEArec template idea (Ricker waveform, point-source spatial decay) but emits data one chunk at a
time as a bounded-memory generator, never materializing the full recording on disk.

\textbf{Rust in scientific computing.} The PyO3 ecosystem \citep{PyO3} lets Rust libraries expose
Python APIs while releasing the GIL during computation, pairing C-level speed with Python usability.
Rayon \citep{Rayon} supplies data-parallel iterators over chunks through a work-stealing thread pool,
with no subprocess overhead and no pickle-based inter-process communication.

\section{Architecture}

\subsection{The \texttt{ChunkSource} trait}

Segovia's central abstraction is the \texttt{ChunkSource} trait, an iterator over fixed-size
$(\text{samples} \times \text{channels})$ \texttt{i16} buffers. Each \texttt{next()} call returns
exactly $\text{chunk\_samples} \times n_{\text{channels}}$ samples in row-major (samples $\times$
channels) order, the layout downstream spike sorters expect. Three production implementations ship:

\begin{itemize}
  \item \textbf{\texttt{SpikeGlxReader}} reads SpikeGLX \texttt{.bin} + \texttt{.meta} pairs through
  memory-mapped I/O. Chunks are sliced from the memory map with zero copy and no decompression.
  \item \textbf{\texttt{ZarrReader}} reads chunked Zarr arrays (gzip, zstd, Blosc codecs) through the
  \texttt{zarrs} crate, aligning Zarr chunk boundaries to the requested chunk size where possible.
  \item \textbf{\texttt{CbinReader}} reads mtscomp-compressed IBL \texttt{.cbin} recordings. Each chunk
  is decompressed on its own through a positioned read with \texttt{flate2} (zlib). This is the format
  used for the real-data benchmark in Section~\ref{sec:real}.
\end{itemize}

\subsection{The preprocessing chain}

A \texttt{ChunkSource} is consumed by \texttt{preprocess(chunk\_source, config)}, which chains three
operations.

\textbf{Bandpass filter.} A 5th-order Butterworth filter with zero-phase response
(\texttt{sosfiltfilt}) is applied to each chunk. Zero-phase filtering needs a look-ahead
(\texttt{margin}) of future samples for the backward pass, so Segovia fetches \texttt{margin} samples
past the current chunk boundary and drops them from the output, keeping the output time-aligned. This
adds a bounded look-ahead latency of $\text{margin} / f_s$ (50\,ms at the default
$\text{margin} = 1500$ samples, $f_s = 30\,000$\,Hz). A causal single-pass mode that removes this
look-ahead is planned but not yet built.

\textbf{Common median reference (CMR).} The median across all channels is subtracted sample by sample.
This suppresses motion artefacts and common-mode electrical noise without a reference electrode. It
runs after bandpass filtering on each chunk on its own.

\textbf{Global ZCA whitening.} The whitening matrix is estimated from the first \texttt{n\_calib}
samples (default 60\,000; 2\,s at 30\,kHz), then stored and reused for every later chunk. Whitening
runs in \texttt{f32} to avoid integer overflow, and the output is \texttt{f32}.

\textbf{Parallelism and GIL release.} The chain is Rayon-parallelized across time-chunks. When
$\text{batch} > 1$, several chunks run concurrently through Rayon's work-stealing pool. Cross-chunk
filter state (the IIR initial conditions for the forward and backward passes) is carried
deterministically: the state for chunk $k$ is seeded from the last \texttt{margin} samples of chunk
$k-1$, so the output is identical for any thread count. The Python GIL is released through PyO3's
\texttt{allow\_threads} for the whole Rust computation, so Python threads blocked on I/O or other work
are not held.

\textbf{Memory bound.} Peak memory is bounded by
\[
M = \text{batch} \times (\text{chunk\_samples} + 2 \times \text{margin}) \times n_{\text{channels}}
\times \text{sizeof}(\texttt{f32}),
\]
which does not depend on recording length. At $\text{batch} = 1$, $\text{chunk} = 9000$,
$\text{margin} = 1500$, $n_{\text{channels}} = 384$, the bound is about 0.08\,GB. Observed peak RSS on
real data runs higher (0.18--0.49\,GB) because of the Python runtime, the decompressor, and the
whitening matrix, but it still scales only with chunk size, not with recording length.

\subsection{Built-in streaming simulator}

\texttt{SyntheticEphysReader} is a \texttt{ChunkSource} that emits arbitrarily long synthetic
multi-channel ephys streams without touching a file. Each chunk is generated on demand from this
model:

\begin{itemize}
  \item \textbf{Spike templates.} Each unit has a Ricker (Mexican-hat) temporal waveform of amplitude
  $A$ convolved with a point-source spatial decay $V(r) = A \times d_{\perp} / r$, where $r$ is the
  probe-to-source distance and $d_{\perp}$ is the perpendicular component. Waveforms are precomputed at
  initialisation and stamped into the chunk buffer at spike times.
  \item \textbf{Spike times.} Each unit fires on its own as a Poisson process with rate $\lambda$.
  \item \textbf{Noise.} Additive white Gaussian noise of standard deviation $\sigma$ is added to all
  channels.
  \item \textbf{PRNG.} A SplitMix64 seed generator feeds per-unit xoshiro256++ generators. The PRNG is
  deterministic across platforms and chunk sizes: split the recording into different chunk sizes and
  the output is bit-identical.
\end{itemize}

The simulator's footprint is one chunk buffer plus the (small, bounded) template bank. It serves two
roles: a standalone benchmark source for systems metrics (Section~\ref{sec:synth}), and the basis for
\texttt{ground\_truth()}, which returns $(\text{sample\_index}, \text{unit\_id},
\text{peak\_channel})$ tuples for MEArec-style spike-sorter accuracy evaluation.

\section{Evaluation methodology}

We follow the replay-at-acquisition-rate method: prerecorded or synthetic data are streamed from disk,
or generated, at the true 30\,kHz sampling rate, and per-chunk end-to-end latency is measured with
nanosecond-precision monotonic clocks (\texttt{perf\_counter\_ns}). A chunk meets its real-time
deadline when its compute latency is at most the chunk period ($\text{chunk\_samples} / f_s$). Deadline
adherence is the fraction of chunks satisfying that bound. The first three chunks of each run are
dropped as warm-up.

\textbf{Metrics reported:} per-chunk latency mean, standard deviation (jitter), p95, p99, and max;
deadline adherence; peak whole-process RSS (sampled every 20\,ms in-process); and sustained throughput
(total output bytes divided by total elapsed wall time).

\textbf{SpikeInterface baseline.} SI runs in a separate venv (\texttt{.venv-si},
\texttt{spikeinterface==0.102.3}) in the same process, driven by sequential
\texttt{get\_traces(start\_frame, end\_frame)} calls with $\text{n\_jobs} = 1$, the online analog of
Segovia's $\text{batch} = 1$. The filter margin is matched (\texttt{bandpass\_filter(margin\_ms=50.0)}).
Whitening uses \texttt{mode="global"} with random calibration chunks (the SI default); Segovia uses the
first 60\,000 samples. Both differences touch only warm-up and are excluded by the three-chunk discard.

\textbf{Machine:} Windows, 8 physical / 16 logical cores, 7.8\,GB RAM.

\section{Results}

\subsection{Real IBL AP-band recording}
\label{sec:real}

Source: \texttt{\_spikeglx\_ephysData\_g0\_t0.imec0.ap.cbin} from the IBL reproducible ephys dataset
\citep{IBL2025} (mtscomp-compressed, 385 channels including sync, 30\,kHz).

\textbf{Full length, 300\,ms budget (steady state, 55.8\,min, 11{,}167 chunks).} This is the
representative sustained-streaming measurement and the headline online result (Table~\ref{tab:steady}).

\begin{table}[h]
\centering
\small
\caption{Full-length steady-state streaming, 300\,ms chunk budget, real IBL recording.}
\label{tab:steady}
\begin{tabular}{lrrrrrrrr}
\toprule
Engine & Mean & p95 & p99 & Max & Jitter & Adher. & Peak & Thpt \\
 & (ms) & (ms) & (ms) & (ms) & (ms) & & RSS & (MB/s) \\
\midrule
Segovia & 179.2 & 256.3 & 277.0 & \textbf{334.5} & \textbf{38.6} & \textbf{99.7\%} & \textbf{0.21\,GB} & 38.1 \\
SI online & 205.3 & 301.4 & 355.0 & 932.0 & 60.5 & 94.7\% & 0.41\,GB & 33.3 \\
\bottomrule
\end{tabular}
\end{table}

Segovia leads on every axis. The largest separations are peak memory (2$\times$), maximum latency
(334 vs 932\,ms), p99 (277 vs 355\,ms), and jitter (38.6 vs 60.5\,ms). The deadline-adherence lead is
real but narrow at steady state, 99.7\% against 94.7\%.

\textbf{Cold-start, first 60\,s (1{,}800{,}000 samples).} The per-chunk sweep in
Table~\ref{tab:cold} covers the cold window: first reads, cold page cache, and library warm-up beyond
the three-chunk discard, where SpikeInterface's per-chunk \texttt{get\_traces} setup cost is highest
and the adherence gap is widest. That gap narrows to the steady-state figures above over the full
recording, so these numbers describe cold-start rather than sustained operation. Only the 300\,ms leg
has a full-length counterpart in Table~\ref{tab:steady}; the 100\,ms and 1000\,ms rows here remain
60\,s cold-start measurements.

\begin{table}[h]
\centering
\small
\caption{Cold-start per-chunk sweep, first 60\,s of the real IBL recording.
597 / 197 / 57 chunks measured per row.}
\label{tab:cold}
\begin{tabular}{llrrrrrrrr}
\toprule
Chunk & Engine & Mean & p95 & p99 & Max & Jitter & Adher. & Peak & Thpt \\
 & & (ms) & (ms) & (ms) & (ms) & (ms) & & RSS & (MB/s) \\
\midrule
100\,ms & Segovia & 92.9 & 118.4 & 122.0 & 127.9 & 15.4 & \textbf{74.2\%} & \textbf{0.21\,GB} & 24.7 \\
100\,ms & SI online & 112.0 & 250.0 & 275.2 & 302.8 & 48.5 & 64.2\% & 0.46\,GB & 20.5 \\
300\,ms & Segovia & 194.5 & 250.1 & 256.4 & 292.5 & 34.7 & \textbf{100\%} & \textbf{0.28\,GB} & 35.2 \\
300\,ms & SI online & 245.8 & 355.7 & 365.7 & 407.1 & 67.8 & 69.5\% & 0.52\,GB & 27.9 \\
1000\,ms & Segovia & 617.3 & 692.3 & 705.9 & 707.9 & 77.5 & \textbf{100\%} & \textbf{0.49\,GB} & 37.4 \\
1000\,ms & SI online & 786.0 & 831.3 & 947.5 & 1036.6 & 42.9 & 98.2\% & 0.74\,GB & 29.1 \\
\bottomrule
\end{tabular}
\end{table}

In this cold window at 300\,ms, Segovia reaches 100\% adherence at 0.28\,GB while SpikeInterface online
reaches 69.5\% at 0.52\,GB, a wider gap than the 99.7\% vs 94.7\% at steady state, because SI pays its
per-chunk \texttt{get\_traces} setup and margin re-decode most heavily before the run warms up (each
call re-reads and re-decodes the 50\,ms filter-margin neighbourhood, with no cross-chunk pipelining;
SI cold p99 is 366\,ms against Segovia's 256\,ms). At 100\,ms both engines fall short of 100\%
adherence because the mtscomp zlib decode is memory-bandwidth bound (established in the project's ADR
0013): the Rust compute meets the deadline but the decode does not. Peak RSS tracks only chunk size on
real data, matching the analytical bound, and Segovia's RSS sits below SI's at every chunk size.

\subsection{Synthetic recordings (Segovia-only baseline)}
\label{sec:synth}

Source: \texttt{SyntheticEphysReader}, 384 channels, 60\,s, 30\,kHz, 20 units, 5\,Hz Poisson firing,
10\,\textmu V noise, seed 0 (Table~\ref{tab:synth}).

\begin{table}[h]
\centering
\small
\caption{Synthetic uncompressed stream, Segovia only.}
\label{tab:synth}
\begin{tabular}{lrrrrrrrr}
\toprule
Chunk & Mean & p95 & p99 & Max & Jitter & Adher. & Peak & Thpt \\
 & (ms) & (ms) & (ms) & (ms) & (ms) & & RSS & (MB/s) \\
\midrule
100\,ms & 61.3 & 68.4 & 73.7 & 83.1 & 3.6 & \textbf{100\%} & 0.15\,GB & 37.2 \\
300\,ms & 147.9 & 163.9 & 167.4 & 178.1 & 8.6 & \textbf{100\%} & 0.23\,GB & 45.5 \\
1000\,ms & 617.8 & 650.9 & 664.9 & 672.8 & 18.6 & \textbf{100\%} & 0.50\,GB & 36.2 \\
\bottomrule
\end{tabular}
\end{table}

On synthetic uncompressed streams Segovia holds 100\% deadline adherence at every chunk size, with much
lower jitter than on real compressed data (3.6\,ms against 15.4\,ms at 100\,ms chunks). That lower
jitter comes from the absence of zlib decode variance. Throughput clears 22\,MB/s in every
configuration.

\subsection{Batch throughput}

The online comparison above is separate from batch throughput, where SpikeInterface's parallel
\texttt{ChunkRecordingExecutor} ($\text{n\_jobs} = N$) runs its C/MKL kernels across all cores, the
regime SI was designed for. We reran this on the full 55.8-minute (29\,GB compressed) real IBL
recording, with Segovia at a pinned batch of 4 (the throughput/memory optimum on this host, ADR 0017;
the footprint is about $0.17\,\text{GB} \times \text{batch} + 0.5\,\text{GB}$), in Table~\ref{tab:batch}.

\begin{table}[h]
\centering
\small
\caption{Full-scale batch throughput, full 55.8-minute real IBL recording.}
\label{tab:batch}
\begin{tabular}{lrr}
\toprule
Engine & Wall time & Peak RSS \\
\midrule
Segovia (batch 4) & \textbf{806\,s} & \textbf{1.19\,GB} \\
SpikeInterface (thread pool, n\_jobs = 8) & 923\,s & 2.18\,GB \\
SpikeInterface (process pool, n\_jobs = 8) & 1022\,s & 4.42\,GB \\
\bottomrule
\end{tabular}
\end{table}

At this pinned batch Segovia used less peak memory than both SI executor modes and finished in less
wall time, in a single run. The memory advantage is the robust, deterministic finding, confirmed
file-size-independent at full scale: batch-1 peak moved only 0.9\% from a ten-minute slice to the full
recording. The wall-time margin is a single-machine, single-run measurement and should be replicated
with confidence intervals across machines before being read as a general speed claim. An earlier report
that Segovia only ties SI in batch (ADR 0013) was taken at the unpinned auto default, which
oversubscribes memory-bandwidth-bound work on many-core hosts (ADR 0017). We verified output
equivalence: without whitening, Segovia matches SI to 0.0035\% (bandpass + CMR), and the whitened-run
difference comes from SI's random-subset whitening calibration against Segovia's deterministic
first-60{,}000-sample calibration, not from a computational error.

\section{Discussion}

\textbf{Online versus batch.} The results draw a clear line: batch tools driven one chunk at a time pay
a per-chunk overhead that a purpose-built streaming engine avoids. This follows from SpikeInterface's
design goal rather than any flaw in it. Closed-loop applications that need sub-300\,ms preprocessing
latency with bounded memory should reach for a streaming-first tool; batch tools stay the right choice
for offline analysis.

\textbf{Non-causal filter and look-ahead.} The zero-phase Butterworth filter introduces a 50\,ms
look-ahead (the $\text{margin} / f_s$ term), so true end-to-end latency from sample acquisition to
preprocessed output is $\text{chunk\_period} + 50\,\text{ms} + \text{compute\_latency}$, not compute
latency alone. At 300\,ms chunks that is about 400--450\,ms total. A causal single-pass mode would
remove the 50\,ms look-ahead at the cost of phase distortion, and is future work. Every latency figure
in this paper is compute latency; the look-ahead is reported separately by the harness.

\textbf{Synthetic-data limits.} The synthetic simulator does not reproduce the exact noise statistics
of real recordings, a documented limit of the MEArec-style approach. The metrics reported here
(latency, memory, deadline adherence) are systems metrics that depend on data shape and scale rather
than biological fidelity, so the synthetic results support the claims made. We keep the real IBL run
for external validity, and it exposes the real decoding bottleneck (mtscomp zlib) that synthetic
uncompressed data does not.

\textbf{Single-machine measurements.} All benchmarks run on one machine (Windows, 8 physical cores,
7.8\,GB RAM), so performance on other hardware will differ. The analytical memory bound
$M = \text{batch} \times (\text{chunk} + 2 \times \text{margin}) \times \text{channels} \times
\text{sizeof}(\texttt{f32})$ is machine-independent; the latency and throughput figures are not.

\textbf{Future work.} A causal single-pass filter mode to remove the 50\,ms look-ahead, and
parallelized mtscomp decompression to lift the 100\,ms budget limit on compressed data.

\section{Conclusion}

Segovia is a streaming, bounded-memory preprocessing engine for Neuropixels-scale electrophysiology
that satisfies two hard constraints for near-real-time work, real-time deadline adherence and
file-size-independent memory, which batch tools do not offer in the online regime. Streaming the full
55.8-minute real IBL recording at a 300\,ms chunk budget, Segovia meets its deadline on 99.7\% of
chunks at 0.21\,GB against 94.7\% at 0.41\,GB for SpikeInterface online, and it holds larger margins in
peak memory, maximum latency, and jitter. The engine ships as \texttt{pip install segovia} (Python,
PyPI) and \texttt{cargo add segovia} (Rust, crates.io) under AGPL-3.0-or-later.

\section*{Acknowledgements}

The name Segovia honors Claudio Segovia (1983--2009), a friend who died of leukemia at 26.

\bibliographystyle{plainnat}
\bibliography{paper}

\end{document}
